\chapter{Homework 2 (due 9/11)}


\subsection{Pauli matrices}

(1) Diagonalize the 3 Pauli matrices and find out their eigenvalues and eigenkets:
\begin{align}
\sigma_x=\bpm0&1\\1&0\epm,~~\sigma_y=\bpm0&-\imth\\ \imth&0\epm,~~\sigma_z=\bpm1&0\\0&-1\epm
\end{align}

(2) For an arbitrary function $f(\cdot)$ from complex numbers to complex numbers, if $\vec n$ is a normalized vector in three dimensions and $\theta$ is a real number, prove the following relation
\begin{align}
f(\theta\vec n\cdot\vec\sigma)=\frac{f(\theta)+f(-\theta)}2\hat 1+\frac{f(\theta)-f(-\theta)}2\vec n\cdot\vec\sigma
\end{align}
where $\vec\sigma=(\sigma_x,\sigma_y,\sigma_z)$ are three Pauli matrices.

Hint: a function $f(\hat M)$ of a matrix (operator) $\hat M$ can be defined by the Taylor expansion:
\begin{align}
f(\hat M)=\sum_{n=0}^\infty\frac{\hat M^n}{n!}\frac{\text{d}^{n}f}{\text{d}x^n}|_{x=0}
\end{align}



\subsection{Linear algebra}

(1) Use the Cauchy-Schwarz inequality to prove the triangle inequality
\begin{align}
|\dket{u}+\dket{v}|\leq|\dket{u}|+|\dket{v}|,~~~\forall~\dket{u},\dket{v}\in V
\end{align}

(2) A linear operator $\hat P$ is a projection operator (or a projector) if $\hat P^2=\hat P$. Prove that the eigenvalues of a projector are all either $0$ or $1$. 

(3) \textbf{Bonus} A projector has the following form:
\begin{align}
\hat P=\alpha\dket{1}\dbra{1}+\beta\dket{2}\dbra{2}
\end{align}
where $\dket{1}$ and $\dket{2}$ are two normalized vectors in the Hilbert space that are linearly independent. What conditions should coefficients $\alpha$ and $\beta$ satisfy in order for $\hat P$ defined above to be a projector (i.e. to satisfy $\hat P^2=\hat P$)?

%(4) Prove that a positive operator $\hat O$ satisfying 
%\begin{align}
%\expval{v|\hat O|v}\geq0,~~~\forall~\dket{v}\in\mathbb{V}
%\end{align}
%is necessarily Hermitian.

%Hint: any linear operator can be written as $O=H_1+\imth H_2$, where $H_{1,2}$ are two Hermitian operators.t

%(4) \textbf{ Bonus} Prove that a Hermitian linear operator $\hat O=\hat O^\dagger$ can be diagonalized, referring to how the spectral decomposition theorem is proved in Box 2.2 of \nielsen.


(4) Assume Hamiltonian $H$ commutes with two linear operators, $A_1$ and $A_2$, which do not commute with each other:
\begin{align}
[H,A_1]=[H,A_2]=0,~~~[A_1,A_2]\neq0
\end{align}
Show that the eigenstates of $H$ are in general degenerate. 

(5) Assume two observables $A_1$ and $A_2$ anticommutes with each other, i.e.
\begin{align}
    \{A_1,A_2\}=A_1A_2+A_2A_1=0
\end{align}
Show that for any eigenstate $\dket{\psi_1}$ of observable $A_1$ with a nonzero eigenvalue, its expectation value w.r.t. $A_2$ must vanish
\begin{align}
    \expval{\psi_1|A_2|\psi_1}=0
\end{align}
Comment: One such example is Pauli matrices $A_1=\sigma_x$ and $A_2=\sigma_y$. 


(5) \textbf{Bonus} For two linear operators satisfying $[\hat A,\hat B]=\alpha\hat 1$ where $\hat 1$ is the identity operator, prove the following identity:
\begin{align}
[\hat A,e^{\hat B}]=\alpha e^{\hat B}
\end{align}

Hint: first show that
\begin{align}
\hat f(x)=e^{-x \hat B}\hat Ae^{x\hat B}=\hat A+x\alpha\hat 1
\end{align}

\subsection{Consecutive measurements of incompatible operators}

Two consecutive projective measurements, for Pauli matrices $\hat \sigma_x$ and $\hat \sigma_z$ respectively, are carried out one after another on a quantum state of a single qubit (i.e. a Hilbert space $\mathbb{V}=\mathbb{C}^2$). In the first procedure, we first measured $\hat\sigma_x$ with an outcome $+1$ and probability $p(X=+1)$, and then measured $\hat\sigma_z$ with an outcome $+1$ with probability $p(Z=+1|X=+1)$. In the second procedure, we switched the order of the two measurements, i.e. we first measured $\hat\sigma_z$ with an outcome $+1$ and probability $p(Z=+1)$, and then measured $\hat\sigma_x$ with an outcome $+1$ with probability $p(X=+1|Z=+1)$. It is found that the probabilites do no depend on the order of the two incompatible measurements, i.e.
\begin{align}\label{hw3.1}
p(X=+1)=p(X=+1|Z=+1),~~~p(Z=+1|X=+1)=p(Z=+1)
\end{align}

(1) Find a single-qubit pure state $\dket{\psi}$ that satisfies the condition (\ref{hw3.1}). 

(2) If the quantum state is a single-qubit mixed state, it is generally described by a density matrix of the following form:
\begin{align}
\hat\rho=\frac{\hat 1+\vec r\cdot\vec\sigma}2,~~~\vec r=(x,y,z)\in\mathbb{R}^3,~~~|\vec r|\leq1.
\end{align}
Find those mixed states that satisfy the same condition (\ref{hw3.1}). 

\subsection{Quantum Zeno effect}

Consider a single qubit evolving under the following Hamiltonian:
\begin{align}
\hat H=-\mu_BB\hat\sigma_z
\end{align}
where $\mu_B=\hbar e/2m_e$ is the Bohr magneton. We define time $T$ as follows:
\begin{align}
T=\frac{\pi m_e}{eB}
\end{align}
Starting from the initial state
\begin{align}
\dket{\psi_0}=\dket{+}=\frac{\dket{\uparrow}+\dket{\downarrow}}{\sqrt2}
\end{align}
at time $t=0$, we evolve the system by time $T/n$, followed by a measurement of $\hat\sigma_x$, and then repeat the same 2-step (evolution by $\frac{T}{n}$ + $\hat\sigma_x$ measurement) procedure $n$ times. What is the probability at time $T$ (after $n$ such 2-step procedures) to produce the measurement outcome of $\sigma_x=+1$?